\chapter{Reproducing the Results}
\label{app:repro}

The numerical results of Chapter~\ref{ch:case} are produced in two steps from the film file. The first step decodes the film and takes the longest; the second runs in seconds.

\begin{lstlisting}[language=bash,numbers=none]
# 1. Frame features, detector events, audition clips
python3 code/analyze-soundtrack-v3.1.py drive-through-fire-2026.mkv \
    --title "Drive Through Fire" --out sound-analysis-v3.1

#    (or, reusing an existing version 3 frame table)
python3 code/analyze-soundtrack-v3.1.py drive-through-fire-2026.mkv \
    --from-frames sound-analysis-v3/frames.csv --out sound-analysis-v3.1

# 2. Dialogue-referenced measures, figures, and LaTeX macros
python3 code/occupation.py \
    --frames sound-analysis-v3.1/frames.csv \
    --events sound-analysis-v3.1/events.csv \
    --film drive-through-fire-2026.mkv \
    --subtitles drive-through-fire-2026.srt \
    --out results

#    supplementary statistics and the segment figure
python3 code/occupation-extras.py --results results

# 3. Blinded annotation set with random controls (Chapter 9)
python3 code/blind-audition.py drive-through-fire-2026.mkv \
    --audition sound-analysis-v3.1/audition \
    --frames sound-analysis-v3.1/frames.csv --out blind --controls 40

# 4. Build the manuscript
pdflatex main.tex && pdflatex main.tex
\end{lstlisting}

Step~2 decodes the film once more to measure BS.1770 loudness (cached in \texttt{results/loudness.csv} for later runs), estimates the subtitle offset, and writes \texttt{results/results.tex}, which the manuscript reads at compile time, together with \texttt{fig-occupation.png}, \texttt{fig-exceedance.png}, \texttt{fig-blocks.png}, a per-second table \texttt{seconds.csv}, and \texttt{occupation.json}. If \texttt{results/results.tex} is absent, every dependent value prints as \pending.

Dependencies are \texttt{ffmpeg}, Python~3 with \texttt{numpy}, \texttt{scipy}, \texttt{pandas}, \texttt{matplotlib}, \texttt{soundfile}, and \texttt{librosa}, and a standard \LaTeX{} installation.

\chapter{Annotation Instrument}
\label{app:instrument}

\section{Label vocabulary}

Source labels, of which one or more may be assigned to a clip, joined with \texttt{+}:

\begin{center}\small
\begin{tabular}{llll}
DRUM & MUSIC\_SWELL & MUSIC\_OTHER & SIREN \\
ALARM & GUNSHOT & EXPLOSION & PUNCH \\
CRASH & GLASS & ENGINE & TIRES \\
YELLING & SCREAM & CROWD & DIALOGUE \\
FALSE\_POSITIVE & UNSURE & & \\
\end{tabular}
\end{center}

\section{Fields}

\begin{center}\small
\begin{tabularx}{\linewidth}{lX}
\toprule
Field & Content \\
\midrule
\texttt{clip} & Random identifier assigned at blinding \\
\texttt{label} & Source labels from the vocabulary \\
\texttt{role} & voice, music, effect, ambience, or mixed \\
\texttt{intensity} & 1 (barely noticeable) to 5 (overwhelming) \\
\texttt{would\_reduce} & yes or no: would the listener reduce the volume \\
\texttt{notes} & Free text \\
\bottomrule
\end{tabularx}
\end{center}

Detector category, score, film time, and whether the clip is a random control are held in a separate key file (\texttt{key.csv}, written by \texttt{blind-audition.py}) and joined to the labelled sheet on the \texttt{clip} column only after labelling is complete. Random control clips appear in the key with category \texttt{RANDOM\_CONTROL}.

\chapter{Code}
\label{app:code}

The four programs used in this study are reproduced in full. All are released into the public domain.

\section{Feature extraction and detection}
\lstinputlisting[language=Python]{code/analyze-soundtrack-v3.1.py}

\section{Dialogue-referenced occupation analysis}
\lstinputlisting[language=Python]{code/occupation.py}

\section{Supplementary statistics}
\lstinputlisting[language=Python]{code/occupation-extras.py}

\section{Blinded annotation set}
\lstinputlisting[language=Python]{code/blind-audition.py}

\section{Volume behaviour logger for \texttt{mpv}}
\lstinputlisting{code/throttle-log.lua}
